\documentclass{SibFU-docs}

% доп. пакеты
\usepackage{graphicx}
\usepackage{float}
\usepackage{hyperref}
\usepackage{caption}
\usepackage{subcaption}
\usepackage{amsmath}
% предотвартим ошибку с \Bbbk перед загрузкой пакета amssymb
\let\Bbbk\relax
\usepackage{amssymb}
\usepackage{booktabs}
\usepackage{tabularx}
\usepackage{longtable}
\usepackage{enumitem}
\usepackage{multirow}

% гиперссылки
\hypersetup{
    colorlinks=true,
    linkcolor=black,
    citecolor=blue,
    urlcolor=blue
}


% библиография
\addbibresource{report.bib}

% титульный лист
\begin{document}
\setlist[itemize]{left=1.3cm, itemsep=1pt, topsep=2pt} % для выравнивания списков
\setlist[enumerate]{left=1.3cm, itemsep=1pt, topsep=2pt}
\renewcommand{\contentsname}{}
\mycovertitle
{Гуманитарный институт}
{Кафедра прикладной информатики в искусстве и интерактивном медиа}
{Сетевая архитектура и базовая безопасность }
{ст. преподаватель И.~Р.~Нигматуллин}
{Захаров И.М., Шикайкова К. С., Тугаринова А. И., Пачковская А. В.,\\
Арсенян А. В., Тетерина П. А., Дегтярева В. В.}
 
\begin{center}
\bfseries РЕФЕРАТ
\end{center}
\noindent

Реферат по теме: «Сетевая архитектура и базовая безопасность » содержит 41 страницу, 7 таблиц и 24 использованных источников.

Ключевые слова: АУТЕНТИФИКАЦИЯ, АВТОРИЗАЦИЯ, ИДЕНТИФИКАЦИЯ, МНОГОФАКТОРНАЯ АУТЕНТИФИКАЦИЯ, КОНТРОЛЬ ДОСТУПА, МЕЖСЕТЕВЫЕ ЭКРАНЫ, IP-АДРЕСАЦИЯ, МАРШРУТИЗАЦИЯ, ЛОКАЛЬНЫЕ И ГЛОБАЛЬНЫЕ СЕТИ, ВРЕДОНОСНОЕ ПО, РЕЗЕРВНОЕ КОПИРОВАНИЕ.

В работе исследуются теоретические и практические аспекты обеспечения безопасного доступа к ресурсам компьютерных сетей. Рассматриваются основные принципы построения сетей и их архитектур, механизмы IP-адресации и маршрутизации, особенности организации локальных и глобальных сетей. Отдельное внимание уделяется функциям межсетевых экранов в защите периметра, методам противодействия вредоносному программному обеспечению, стратегиям резервного копирования и восстановления данных, а также современным подходам к аутентификации и управлению доступом пользователей.

Основные выводы: надёжная аутентификация и корректно настроенный контроль доступа невозможны без целостного понимания сетевой архитектуры, механизмов маршрутизации и сегментации, использования межсетевых экранов и средств защиты конечных точек. Комплексное применение криптографических протоколов, многофакторной аутентификации, антивирусной защиты и регулярного резервного копирования позволяет существенно снизить риски компрометации учётных записей и потери данных.

Рекомендации: при проектировании и эксплуатации информационных систем следует реализовывать многоуровневую защиту (defense in depth), сочетая надёжную аутентификацию, жёсткие политики контроля доступа, сетевую фильтрацию, средства защиты от вредоносного ПО и формализованные процедуры резервного копирования. Необходимо регулярно проводить аудит настроек доступа, обновлять программное обеспечение и обучать пользователей основам безопасной работы в сети.


\newpage

\thispagestyle{empty}
\begin{center}
\bfseries СОДЕРЖАНИЕ
\end{center}
\vspace{-6pt}
\tableofcontents
\clearpage

% фантомные боли для секций, чтобы не было нумераций кроме тела реферата
\section*{}
\vspace*{-2.5em}
\begin{center}\textbf{ВВЕДЕНИЕ}\end{center}
\phantomsection
\addcontentsline{toc}{section}{Введение}

Компьютерные сети стали неотъемлемой частью современного информационного общества, обеспечивая обмен данными, совместное использование ресурсов и глобальную коммуникацию. От локальных сетей внутри организаций до всемирной сети Интернет --- компьютерные сети формируют фундамент цифровой экономики и электронного взаимодействия.

Однако с ростом значимости сетевой инфраструктуры возрастают и угрозы информационной безопасности. Несанкционированный доступ, перехват данных, вредоносное программное обеспечение и сетевые атаки представляют серьёзную угрозу для пользователей и организаций. В этих условиях механизмы аутентификации и контроля доступа становятся критически важными компонентами защиты информационных ресурсов.

Современное состояние проблемы характеризуется многоуровневым подходом к безопасности. Классические механизмы парольной аутентификации дополняются многофакторной верификацией, биометрическими технологиями и криптографическими протоколами. Сетевая безопасность строится на основе эшелонированной защиты, включающей межсетевые экраны, системы обнаружения вторжений, антивирусные средства и инструменты резервного копирования.

Актуальность темы определяется возрастающей зависимостью организаций и частных лиц от сетевых сервисов, ростом киберпреступности и усложнением ландшафта угроз. По данным международных исследований, ущерб от кибератак исчисляется триллионами рублей ежегодно, а средний срок обнаружения вторжения в корпоративные сети может составлять месяцы. В условиях удалённой работы и цифровой трансформации обеспечение аутентификации и доступа становится задачей первостепенной важности.

Новизна работы заключается в комплексном рассмотрении аутентификации и контроля доступа в контексте всего стека сетевых технологий --- от физической топологии и протоколов маршрутизации до прикладных механизмов защиты конечных точек и процедур восстановления данных.

Цель работы: изучение принципов построения компьютерных сетей и механизмов обеспечения безопасного доступа к информационным ресурсам; выявление взаимосвязи между сетевой архитектурой, механизмами аутентификации и средствами защиты от угроз.

Задачи: рассмотреть основные принципы построения компьютерных сетей и их архитектур; изучить механизмы IP-адресации и маршрутизации данных; проанализировать особенности локальных и глобальных сетей; исследовать роль межсетевых экранов в защите сетевой инфраструктуры; изучить методы защиты от вредоносного программного обеспечения; рассмотреть принципы резервного копирования и восстановления данных как компонента доступности информационных ресурсов.

Методы: системный анализ сетевых архитектур, сравнительный анализ технологий и протоколов, анализ литературы и международных стандартов (ISO/IEC 27001, NIST), изучение практических реализаций систем аутентификации и контроля доступа.

\newpage
% вручную увеличить номер секции и записать в оглавление с номером,
% чтобы в теле не печатался автоматический номер, но в tableofcontents он остался
\refstepcounter{section}
\addcontentsline{toc}{section}{\protect\numberline{\thesection} Основные принципы построения компьютерных сетей и их архитектуры}
\vspace*{-2.5em}
\begin{center}\textbf{\thesection\ ОСНОВНЫЕ ПРИНЦИПЫ ПОСТРОЕНИЯ КОМПЬЮТЕРНЫХ СЕТЕЙ И ИХ АРХИТЕКТУРЫ}\end{center}
\vspace*{-1.5em}
\ManualSubsection{Определение и назначение компьютерных сетей}

Компьютерная сеть представляет собой совокупность компьютеров и других устройств, соединённых между собой каналами связи для обмена данными и совместного использования ресурсов. \cite{tanenbaum2011computer} определяет компьютерную сеть как систему автономных компьютеров, связанных технологией передачи данных. В современном мире компьютерные сети являются основой цифровой коммуникации, охватывая спектр от локальных сетей внутри офиса до глобальной сети Интернет.

Изучение принципов построения компьютерных сетей включает понимание физического размещения устройств через топологии, уровней взаимодействия посредством моделей OSI и TCP/IP, архитектурных решений и используемых протоколов.

\ManualSubsection{Сетевые абоненты и ресурсы}

Сетевые абоненты представляют собой конечные устройства в сети: персональные компьютеры, серверы, мобильные устройства и периферийное оборудование, такое как принтеры и сканеры. Сетевые ресурсы разделяются на аппаратные (принтеры, диски, модемы), программные (файлы, базы данных) и сервисы (электронная почта, веб-сервисы, видеоконференции).

\ManualSubsection{Модели построения сетей}
\vspace*{-1.5em}
\ManualSubsubsection{Одноранговые сети}

В одноранговых сетях (Peer-to-Peer, P2P) все компьютеры равноправны. Каждый компьютер может быть одновременно источником и потребителем ресурсов. Различают неструктурированные P2P-сети, где участники связываются случайным образом, структурированные с чёткой архитектурой взаимосвязей, обеспечивающей эффективный обмен данными, и гибридные, объединяющие элементы обоих подходов для высокой производительности и устойчивости при сохранении децентрализации.

Преимущества одноранговых сетей включают простоту реализации, низкие затраты на оборудование и отсутствие единой точки отказа. К недостаткам относятся сложное управление, низкая производительность при большом количестве компьютеров и слабая централизованная безопасность.

\ManualSubsubsection{Архитектура клиент-сервер}

В архитектуре клиент-сервер выделенные компьютеры выступают в роли серверов (поставщиков ресурсов), а остальные --- в роли клиентов (потребителей ресурсов). \cite{kurose2021computer} подчёркивает, что клиент-серверная модель является фундаментальной для большинства современных сетевых приложений, от HTTP/HTTPS для веб-приложений до TCP/IP для общей сетевой коммуникации.

Преимущества клиент-серверной архитектуры включают централизованное управление, лёгкость обеспечения безопасности, масштабируемость и эффективное использование ресурсов. Недостатками являются высокие затраты на серверное оборудование, зависимость от надёжности сервера и уязвимость к сбоям сервера.

\ManualSubsection{Топологии компьютерных сетей}

Топология сети определяет физическое и логическое расположение устройств и способ их соединения. Существует различие между физической топологией (реальное расположение кабелей) и логической топологией (маршруты передачи сигналов).

\ManualSubsubsection{Физическая топология звезда}

В топологии звезда все устройства подключены к центральному узлу --- коммутатору, маршрутизатору или серверу. Реализуется двусторонняя связь точка-точка между центральным узлом и каждым устройством.

Достоинства топологии звезда включают изоляцию отказов (выход из строя одного компьютера не влияет на работу сети), высокую производительность, лёгкость добавления новых компьютеров и простое обнаружение проблем. Недостатками являются единая точка отказа в виде центрального узла, повышенный расход кабеля и зависимость от производительности центрального устройства. Типичным примером служит локальная сеть в офисе с коммутатором в центре.

\ManualSubsubsection{Физическая топология кольцо}

В топологии кольцо устройства соединены в замкнутый цикл. Сигналы передаются в одном направлении, причём каждый компьютер усиливает сигнал и передаёт дальше. Достоинства включают равномерное распределение нагрузки, хорошую пропускную способность и отсутствие необходимости в центральном узле. Недостатки: выход из строя одного компьютера разрывает кольцо, сложность добавления и удаления устройств, затруднённое диагностирование неисправностей и задержки распространения сигналов.

\ManualSubsubsection{Физическая топология шина}

В топологии шина все устройства подключены к одному общему каналу (магистрали). Сигналы передаются всем узлам одновременно, а на концах кабеля требуются терминаторы. Достоинства: простота реализации, низкие затраты на кабель, изоляция отказов отдельных компьютеров. Недостатки: разрыв кабеля отключает всю сеть, низкая производительность при большом числе узлов, сложность диагностирования проблем и конфликты при одновременной передаче.

\ManualSubsubsection{Древовидная топология}

Древовидная топология представляет собой комбинацию звезды и шины с иерархическим расположением узлов. Один основной узел (корень) управляет подчинёнными уровнями. Достоинства включают масштабируемость, поддержку приоритетов и относительно простое управление. Недостатки: зависимость от узлов более высокого уровня и усложнённая топология по сравнению с простыми схемами. Примером служит сеть из нескольких офисов, соединённых иерархически.

\ManualSubsubsection{Смешанная топология}

На практике часто используется комбинация различных топологий. Например, физическая топология может быть звездой, а логическая --- общей шиной (классический Ethernet), либо комбинация звезды и дерева для расширения сети.

\ManualSubsection{Модель OSI}

Модель Open Systems Interconnection (OSI) была разработана Международной организацией по стандартизации (ISO) для стандартизации сетевого взаимодействия. \cite{stallings2010data} описывает модель OSI как концептуальную основу для понимания сетевых протоколов через семь уровней абстракции.

Модель описывает передачу данных в сетях через семь уровней. Основной принцип построения модели: на каждом уровне добавляется служебная информация (заголовок и хвостовик). Этот процесс называется инкапсуляцией. У отправителя данные проходят сверху вниз, на каждом уровне добавляя заголовки. У получателя происходит деинкапсуляция снизу вверх, последовательно удаляя заголовки.

\vspace{1.5em}

\noindent Таблица 1 – Уровни модели OSI

{\small
\noindent
\begin{tabular}{|p{0.12\textwidth}|p{0.22\textwidth}|p{0.30\textwidth}|p{0.28\textwidth}|}
\hline
\textbf{Уровень} & \textbf{Название} & \textbf{Функции} & \textbf{Примеры} \\
\noalign{\hrule height 1.2pt}
\hline
7 & Прикладной & Сервисы приложений, интерфейс пользователя & HTTP, SMTP, FTP, DNS \\
\hline
6 & Представления & Преобразование формата, шифрование, сжатие & SSL/TLS, JPEG \\
\hline
5 & Сеансовый & Управление сеансом, синхронизация диалога & NetBIOS, RPC \\
\hline
4 & Транспортный & Надёжная/ненадёжная передача, управление потоком & TCP, UDP \\
\hline
3 & Сетевой & Маршрутизация, логическая адресация & IP, ICMP \\
\hline
2 & Канальный & Локальная доставка, управление доступом к среде & Ethernet, Wi-Fi \\
\hline
1 & Физический & Передача физических сигналов & Кабели, сигналы \\
\hline
\end{tabular}
}

\vspace{1em}

Физический уровень отвечает за передачу физических сигналов по сети: электрических сигналов в кабелях, световых импульсов в волоконных кабелях и радиоволн в беспроводных сетях. Канальный уровень обеспечивает доставку данных на соседний узел в локальной сети, управление доступом к среде передачи через MAC, обнаружение и исправление ошибок, упаковку данных в кадры с MAC-адресами.

Сетевой уровень обеспечивает доставку пакетов между узлами в разных сетях, маршрутизацию (выбор пути передачи), логическую адресацию через IP-адреса, фрагментацию и сборку пакетов. Транспортный уровень реализует надёжную или ненадёжную передачу данных, управление потоком, контроль ошибок и мультиплексирование через порты.

Сеансовый уровень управляет установкой и прекращением сеансов, синхронизацией диалога между приложениями и управлением обменом сообщениями. Уровень представления выполняет преобразование формата данных, шифрование и дешифрование, сжатие и распаковку, кодирование символов. Прикладной уровень обеспечивает взаимодействие с приложениями, предоставление сетевых сервисов пользователям и интерфейс между пользователем и сетью.

\ManualSubsection{Модель TCP/IP}

Модель TCP/IP (также называется моделью ARPANET) является практической моделью, на которой построен Интернет. В отличие от концептуальной модели OSI, TCP/IP используется реально и состоит из четырёх уровней.

Уровень сетевого интерфейса объединяет функции физического и канального уровней модели OSI, обеспечивая физическое соединение с сетью, управление доступом к среде и работу с MAC-адресами. Интернет-уровень соответствует сетевому уровню модели OSI, обеспечивая маршрутизацию пакетов, логическую адресацию через IP и управление сетевыми ошибками.

Транспортный уровень соответствует транспортному уровню модели OSI, реализуя надёжную доставку данных через TCP, быструю ненадёжную доставку через UDP, управление портами, сегментацию и сборку данных, идентификацию приложений. Прикладной уровень объединяет функции сеансового, представления и прикладного уровней модели OSI, обеспечивая взаимодействие с приложениями, преобразование формата данных и управление пользовательскими сессиями.

\vspace{1.5em}

\noindent Таблица 2 – Сравнение моделей OSI и TCP/IP

{\small
\noindent
\begin{tabular}{|p{0.22\textwidth}|p{0.35\textwidth}|p{0.35\textwidth}|}
\hline
\textbf{Аспект} & \textbf{OSI} & \textbf{TCP/IP} \\
\noalign{\hrule height 1.2pt}
\hline
Количество уровней & 7 & 4 \\
\hline
Тип модели & Концептуальная & Практическая \\
\hline
Разработка & Теоретическая & Эмпирическая (из реальной практики) \\
\hline
Использование & Обучение, стандартизация & Интернет, реальные сети \\
\hline
Детализация & Высокая & Менее детальная \\
\hline
\end{tabular}
}

\vspace{1em}

Инкапсуляция в TCP/IP происходит следующим образом: приложение формирует данные, транспортный уровень добавляет заголовок TCP с портами и контрольной суммой, интернет-уровень добавляет заголовок IP с IP-адресами, сетевой интерфейс добавляет заголовок Ethernet с MAC-адресами, физический уровень отправляет кадр Ethernet в виде битов. На получающей стороне происходит деинкапсуляция в обратном порядке.

\ManualSubsection{Архитектуры сетей и протоколы}

Ethernet представляет собой наиболее распространённую архитектуру локальных сетей. \cite{spurgeon2014ethernet} описывает эволюцию Ethernet от простой шинной топологии до современных высокоскоростных сетей. Изначально использовалась шинная топология, в настоящее время чаще применяется звезда. Логическая топология остаётся общей шиной. Скорость варьируется от 10 Мбит/с до 100 Гбит/с и выше. Метод доступа CSMA/CD включает прослушивание канала перед передачей, передачу данных при свободном канале, обнаружение коллизий, отправку сигнала задержки при коллизии и повторную попытку передачи.

Протоколы прикладного уровня включают HTTP для передачи веб-страниц, SMTP для отправки электронной почты, POP3 для получения почты, FTP для передачи файлов, DNS для преобразования доменных имён в IP-адреса и SSH для защищённого удалённого доступа. Протоколы транспортного уровня представлены TCP (надёжная доставка с установлением соединения через three-way handshake, контролем потока и управлением перегруженностью) и UDP (быстрая ненадёжная доставка без установления соединения с низкими задержками).

Протоколы сетевого уровня включают IP для логической адресации и маршрутизации пакетов (IPv4 с 32-битными адресами и IPv6 со 128-битными адресами) и ICMP для проверки доступности хоста через ping и отслеживания маршрута через traceroute.

\newpage
\refstepcounter{section}
\addcontentsline{toc}{section}{\protect\numberline{\thesection}IP‑адресация и маршрутизация данных в сетевых окружениях}
\begin{center}\textbf{\thesection\ IP-АДРЕСАЦИЯ И МАРШРУТИЗАЦИЯ ДАННЫХ В СЕТЕВЫХ ОКРУЖЕНИЯХ}\end{center}
\vspace*{-1.5em}
\ManualSubsection{Роль IP-адресации и маршрутизации}

В основе любой современной компьютерной сети, включая глобальный Интернет, лежат два фундаментальных процесса: IP-адресация и маршрутизация. Они обеспечивают точную доставку информации от источника к получателю в сложной распределённой среде, состоящей из миллионов устройств. \cite{comer2018internetworking} определяет IP-адресацию как систему уникальной идентификации каждого устройства в сети, а маршрутизацию --- как процесс определения оптимального пути для передачи данных через последовательность промежуточных сетей и узлов.

Без этих механизмов была бы невозможна работа Интернета, корпоративных сетей, облачных сервисов и Интернета вещей (IoT).

\ManualSubsection{IP-адресация}

IP-адрес (Internet Protocol Address) представляет собой числовой идентификатор, назначаемый сетевому интерфейсу устройства для его распознавания в сети, использующей стек протоколов TCP/IP. Он работает на сетевом уровне (третий уровень модели OSI).

\ManualSubsubsection{Версии IP-протокола}

IPv4 (Internet Protocol version 4) представляет собой 32-битный адрес, записываемый в виде четырёх десятичных чисел от 0 до 255, разделённых точками (например, 192.168.1.1). Обеспечивает около 4,3 миллиарда уникальных адресов. Исчерпание пула адресов привело к разработке технологий для его экономии (NAT, CIDR) и переходу на IPv6.

IPv6 (Internet Protocol version 6) является 128-битным адресом, записываемым в виде восьми групп четырёх шестнадцатеричных цифр (например, 2001:0db8:85a3:0000:0000:8a2e:0370:7334). Обеспечивает практически неограниченное количество адресов (приблизительно \(3.4 \times 10^{38}\)). Имеет встроенные преимущества: упрощённая маршрутизация, автоматическая конфигурация и улучшенная безопасность (IPsec).

Исторически адреса IPv4 делились на классы (A, B, C, D, E). В современной практике используется бесклассовая адресация CIDR (Classless Inter-Domain Routing), позволяющая гибко делить адресное пространство и экономить адреса. \cite{fuller1993cidr} описывает CIDR как механизм агрегации адресов для эффективного использования пространства. Маска подсети определяет, какая часть адреса относится к номеру сети, а какая --- к номеру хоста, записывается как 255.255.255.0 или в формате CIDR как /24 (для адреса 192.168.1.0/24).

\ManualSubsubsection{Типы IP-адресов}

Публичные (глобальные) адреса уникальны в масштабах Интернета и назначаются провайдерами (ISP) и региональными интернет-регистраторами (RIR). Приватные (частные) адреса используются внутри локальных сетей и не маршрутизируются в Интернете. Диапазоны включают 10.0.0.0 -- 10.255.255.255 (класс A), 172.16.0.0 -- 172.31.255.255 (класс B) и 192.168.0.0 -- 192.168.255.255 (класс C).

Специальные адреса включают 127.0.0.1 (localhost, петлевой интерфейс), 0.0.0.0 (адрес по умолчанию/все сети) и 255.255.255.255 (ограниченный широковещательный адрес).

\vspace{1.5em}

\noindent Таблица 3 – Типы IP-адресов

{\small
\noindent
\begin{tabular}{|p{0.20\textwidth}|p{0.38\textwidth}|p{0.36\textwidth}|}
\hline
\textbf{Тип адреса} & \textbf{Характеристика} & \textbf{Примеры / Диапазоны} \\
\noalign{\hrule height 1.2pt}
\hline
Публичные & Уникальны глобально, маршрутизируются в Интернете & Назначаются провайдерами \\
\hline
Приватные & Используются внутри локальных сетей & 10.0.0.0/8, 172.16.0.0/12, 192.168.0.0/16 \\
\hline
Специальные & Служебные адреса & 127.0.0.1 (localhost), 0.0.0.0, 255.255.255.255 \\
\hline
\end{tabular}
}

\vspace{1em}

Статический IP-адрес назначается вручную и не меняется. Динамический IP-адрес автоматически назначается сервером DHCP (Dynamic Host Configuration Protocol) на ограниченное время (аренда). Это стандартный подход для большинства клиентских устройств.

\ManualSubsection{Маршрутизация данных}

Маршрутизация представляет собой процесс выбора пути в одной или нескольких сетях для доставки пакета от источника к пункту назначения. Функцию маршрутизации выполняют маршрутизаторы (роутеры) --- специализированные устройства, анализирующие IP-заголовки пакетов и сверяющие адрес назначения с таблицей маршрутизации.

\ManualSubsubsection{Таблица маршрутизации}

Каждый маршрутизатор и хост имеют таблицу маршрутизации, содержащую записи (маршруты). Сеть назначения представляет целевую IP-сеть (например, 192.168.2.0/24). Маска определяет маску целевой сети. Шлюз (Gateway) указывает IP-адрес следующего маршрутизатора, на который нужно отправить пакет для достижения целевой сети. Интерфейс представляет собой сетевой порт, через который нужно отправить пакет. Метрика отражает числовую стоимость маршрута (чем меньше, тем предпочтительнее), учитывая задержку, пропускную способность и количество хопов (прыжков/переходов).

Алгоритм работы маршрутизатора включает получение IP-пакета, извлечение IP-адреса назначения из заголовка, поиск в таблице маршрутизации наиболее специфичного маршрута (с самой длинной маской), соответствующего адресу назначения, пересылку пакета на указанный в маршруте шлюз через указанный интерфейс. Если подходящий маршрут не найден, используется маршрут по умолчанию (0.0.0.0/0), направляющий пакет к шлюзу провайдера.

\ManualSubsubsection{Типы маршрутизации}

Статическая маршрутизация предполагает, что маршруты добавляются в таблицу вручную администратором. Плюсы включают безопасность, предсказуемость и минимальные накладные расходы. Минусы: не масштабируется, требует ручного изменения при изменениях в топологии сети. Применяется в небольших сетях, для маршрута по умолчанию и в критических сегментах.

Динамическая маршрутизация означает, что маршрутизаторы автоматически обмениваются информацией о сетях с помощью специальных протоколов маршрутизации и самостоятельно строят таблицы. \cite{huitema1995routing} описывает принципы работы динамических протоколов маршрутизации. Плюсы: масштабируемость, адаптивность к изменениям, отказоустойчивость. Минусы: сложность настройки, нагрузка на процессор и каналы связи, потенциальные уязвимости.

Протоколы динамической маршрутизации классифицируются по области действия и алгоритму. IGP (Interior Gateway Protocols) используются для маршрутизации внутри одной автономной системы (AS --- домен под единым управлением). Примерами служат RIP (Routing Information Protocol), простой протокол на основе количества хопов, медленно сходящийся и устаревший; OSPF (Open Shortest Path First), протокол состояния каналов, строящий полную карту сети и вычисляющий кратчайшие пути с помощью алгоритма Дейкстры, обеспечивающий быструю сходимость и масштабируемость; EIGRP (Enhanced Interior Gateway Routing Protocol), проприетарный протокол Cisco, являющийся гибридным протоколом с чертами Distance Vector и Link-State, эффективный и обеспечивающий быструю сходимость.

EGP (Exterior Gateway Protocols) используются для обмена маршрутной информацией между различными автономными системами. BGP (Border Gateway Protocol) является протоколом вектора пути (Path Vector) и основой маршрутизации в Интернете. \cite{rekhter2006bgp} определяет BGP как стандарт междоменной маршрутизации. BGP принимает решения на основе политик (атрибутов AS\_PATH, LOCAL\_PREF, MED), а не только метрики, обеспечивая стабильность и управляемость глобальной маршрутизации.

По алгоритму протоколы разделяются на дистанционно-векторные (RIP, EIGRP), где маршрутизаторы сообщают соседям только свои таблицы маршрутов (вектор расстояний), и протоколы состояния каналов (OSPF, IS-IS), где маршрутизаторы рассылают информацию о состоянии своих каналов всем маршрутизаторам в области, строя общую карту топологии.

\ManualSubsection{Взаимодействие адресации и маршрутизации на практике}

Рассмотрим сценарий: пользователь в домашней сети (192.168.1.5) запрашивает сайт www.example.com (IP 93.184.216.34). Устройство определяет, что адрес 93.184.216.34 не принадлежит локальной сети 192.168.1.0/24. Пакет отправляется на шлюз по умолчанию --- домашний роутер (192.168.1.1). Роутер выполняет трансляцию адресов (NAT), заменяя приватный адрес источника (192.168.1.5) на свой публичный адрес, выданный провайдером (например, 95.165.32.10), и запоминает соответствие для обратного трафика.

Пакет с IP назначения 93.184.216.34 проходит через цепочку маршрутизаторов провайдера. Каждый маршрутизатор использует свою таблицу маршрутизации, построенную по протоколам OSPF (внутри AS) и BGP (между AS), чтобы найти путь к сети, содержащей адрес 93.184.216.34. Пакет достигает сервера example.com, который формирует ответ. Ответный пакет проходит обратный путь, используя те же принципы маршрутизации. Домашний роутер, получив пакет на свой публичный адрес, на основе таблицы NAT переправляет его на внутренний адрес 192.168.1.5.

\newpage
\refstepcounter{section}
\addcontentsline{toc}{section}{\protect\numberline{\thesection}Организация и назначение локальных и глобальных сетей}
\begin{center}\textbf{\thesection\ ОРГАНИЗАЦИЯ И НАЗНАЧЕНИЕ ЛОКАЛЬНЫХ И ГЛОБАЛЬНЫХ СЕТЕЙ}\end{center}
\vspace*{-1.5em}
\ManualSubsection{Локальные вычислительные сети}

Локальная вычислительная сеть (Local Area Network, LAN) представляет собой сеть, ограниченную небольшим географическим пространством. Объединяют относительно небольшое число компьютеров (10--100) в пределах одного помещения (учебный компьютерный класс), здания или учреждения (университета). Раньше использовались в основном для решения вычислительных задач, сейчас в 99\% случаев используются для обмена информацией. Причём от 60\% до 90\% информации циркулирует внутри LAN, не нуждаясь в выходе наружу.

\ManualSubsubsection{Организация LAN}

Коммутаторы (Switch) представляют собой интеллектуальные устройства, анализирующие MAC-адреса и передающие данные только целевому узлу. Различают L2-коммутаторы, работающие с MAC-адресами и поддерживающие VLAN, L3-коммутаторы, работающие с IP-адресами и выполняющие маршрутизацию, и гибридные L2/L3 коммутаторы, объединяющие обе функции и обеспечивающие высокую производительность.

Маршрутизаторы (Router) на границе LAN обеспечивают связь с другими сетями (например, с Интернетом). Точки доступа (Access Point) организуют беспроводное подключение (Wi-Fi) в рамках стандартов IEEE 802.11. Среда передачи включает витую пару (категории 5e/6 и выше), оптическое волокно или радиоволны (Wi-Fi).

\ManualSubsubsection{Назначение локальных сетей}

Основные функции LAN включают совместное использование ресурсов для доступа к общим сетевым принтерам, файловым серверам, сканерам и системам хранения данных (NAS, SAN), высокоскоростную коммуникацию для обеспечения быстрого обмена данными между узлами внутри организации, организацию коллективной работы через корпоративные почтовые серверы, системы группового планирования, мессенджеры и базы данных, а также централизованное управление и безопасность путём администрирования пользователей, политик безопасности и развёртывания ПО с единой точки.

\ManualSubsection{Глобальные вычислительные сети}

Глобальная вычислительная сеть (Wide Area Network, WAN) представляет собой сеть, покрывающую огромные географические области: регионы, страны, континенты. Первая, самая большая и популярная глобальная сеть --- это Интернет. \cite{kleppmann2017designing} рассматривает WAN как основу для распределённых систем и облачных вычислений.

В глобальных сетях существует два режима информационного обмена. Диалоговый режим (режим реального времени), в котором пользователь, получив порцию информации, может немедленно на неё реагировать, называется on-line. В пакетном режиме, называемом off-line, пользователь передаёт порцию информации или принимает её в коротком сеансе связи и на некоторое время отключается от сети.

\ManualSubsubsection{Организация WAN}

Физическая и логическая инфраструктура WAN включает каналы связи, арендуемые у телекоммуникационных провайдеров: выделенные линии (Leased Line) для постоянного соединения точка-точка (T1/E1, T3/E3), цифровые услуги SDH/SONET и MPLS (Multiprotocol Label Switching) --- технологию, обеспечивающую высокоскоростную передачу данных с помощью меток, поддерживающую QoS и виртуальные частные сети, оптоволоконные магистрали с DWDM (Dense Wavelength Division Multiplexing), позволяющие передавать десятки и сотни каналов по одному волокну на разных длинах волн, а также спутниковые и радиоканалы.

Ключевое оборудование и технологии включают маршрутизаторы (Routers) --- основные интеллектуальные устройства WAN, работающие на сетевом уровне (Уровень 3 OSI) и принимающие решения о пересылке пакетов на основе IP-адресов и таблиц маршрутизации. Протоколы маршрутизации для обмена информацией о путях между сетями используют внутренние (IGP) протоколы OSPF (Open Shortest Path First) и EIGRP для маршрутизации внутри одной автономной системы (организации, провайдера), а также внешние (EGP) протоколы BGP (Border Gateway Protocol) --- протокол, «склеивающий» Интернет путём обмена маршрутами между автономными системами (AS), являющийся основным протоколом межсетевого взаимодействия в глобальном масштабе.

Технологии удалённого доступа включают VPN (Virtual Private Network), создающий зашифрованный «туннель» через публичную сеть (Интернет), позволяя удалённым пользователям или филиалам безопасно подключаться к корпоративной LAN, и SD-WAN (Software-Defined WAN) --- современную архитектуру, использующую программное обеспечение для централизованного и гибкого управления подключениями WAN (включая интернет-каналы), автоматически выбирая оптимальный путь для трафика.

\ManualSubsubsection{Назначение глобальных сетей}

Назначение WAN включает интеграцию распределённых ресурсов для объединения удалённых локальных сетей филиалов компании и облачных платформ в единую корпоративную инфраструктуру, обеспечение глобального доступа к информации через функционирование Интернета как всемирной WAN для доступа к веб-сайтам, облачным сервисам (SaaS) и электронной почте, а также глобальную коммуникацию путём поддержки видеоконференцсвязи, международной телефонии (VoIP) и систем удалённого обучения.

\vspace{1.5em}

\noindent Таблица 4 – Сравнительный анализ LAN и WAN

{\small
\noindent
\begin{tabular}{|p{0.22\textwidth}|p{0.36\textwidth}|p{0.36\textwidth}|}
\hline
\textbf{Критерий} & \textbf{Локальная сеть (LAN)} & \textbf{Глобальная сеть (WAN)} \\
\noalign{\hrule height 1.2pt}
\hline
Географический масштаб & Ограниченное помещение, здание, кампус & Неограниченный: город, страна, мир \\
\hline
Скорость передачи & Высокая (100 Мбит/с -- 100 Гбит/с) & Ниже, варьируется (зависит от канала) \\
\hline
Стоимость владения & Относительно низкая & Очень высокая (аренда каналов) \\
\hline
Управление & Единое, централизованное & Распределённое, совместное с провайдерами \\
\hline
Протоколы доступа & Ethernet (IEEE 802.3), Wi-Fi (802.11) & PPP, HDLC, MPLS \\
\hline
Основное назначение & Совместное использование внутренних ресурсов & Связь между удалёнными сетями, глобальный доступ \\
\hline
\end{tabular}
}

\vspace{1em}

Локальные (LAN) и глобальные (WAN) сети представляют собой взаимодополняющие элементы единой сетевой инфраструктуры. LAN служит для эффективного объединения ресурсов внутри ограниченного пространства. WAN выполняет функцию «мостика», связывая разрозненные LAN в глобальное информационное пространство. Понимание принципов организации и назначения каждого типа сетей является базовым для проектирования и администрирования современных информационных систем.

\newpage
\refstepcounter{section}
\addcontentsline{toc}{section}{\protect\numberline{\thesection}Назначение и функции фаерволов в защите сетевой инфраструктуры}
\begin{center}\textbf{\thesection\ НАЗНАЧЕНИЕ И ФУНКЦИИ ФАЕРВОЛОВ В ЗАЩИТЕ СЕТЕВОЙ ИНФРАСТРУКТУРЫ}\end{center}
\vspace*{-1.5em}
\ManualSubsection{Определение и назначение межсетевого экрана}

В современном цифровом мире сеть стала ключевой средой для обмена данными, работы бизнес-приложений и хранения критической информации. Именно поэтому сетевые атаки, эксплуатация уязвимостей и несанкционированный доступ к ресурсам представляют серьёзную угрозу для пользователей и организаций. Файрвол играет роль «первой линии обороны» на границе сети, контролируя сетевой трафик и снижая риск успешных атак.

Файрвол (межсетевой экран, брандмауэр) представляет собой программный или программно-аппаратный компонент, реализующий управляемый барьер между различными зонами доверия, такими как внутренняя сеть и Интернет. \cite{cheswick2003firewalls} определяет межсетевые экраны как критический элемент сетевой безопасности. Его основной задачей является контроль и фильтрация сетевого трафика на основе заранее определённых правил безопасности. Брандмауэр обычно контролирует границу частной сети или хоста. Сетевые экраны обычно устанавливают на отдельных компьютерах, подключённых к сети, или устройствах пользователей и других конечных точках (хостах).

\ManualSubsection{Принцип работы файрволов}

При прохождении трафика брандмауэр должен решить, не опасен ли этот трафик и можно ли его пропустить. То есть он делит трафик на плохой и хороший, доверенный и недоверенный. Фильтрация трафика с помощью брандмауэра осуществляется на основе заранее определённых или динамических правил контроля подключений. Следуя этим правилам, он пропускает или блокирует веб-трафик, проходящий через частную сеть и подключённые к ней устройства.

Для фильтрации трафика все сетевые экраны, независимо от их типа, используют ту или иную комбинацию следующих атрибутов: адрес источника (устройство, пытающееся подключиться к вашей сети), адрес назначения (устройство вашей сети, к которому пытается подключиться внешнее устройство), тип сообщения (то, что пытается передать подключающееся устройство), пакетные протоколы («язык» подключения, который используется для передачи сообщения; протоколы TCP/IP используются хостами для «общения» друг с другом через интернет, а также внутри сети или подсети чаще других) и протоколы прикладного уровня (распространённые протоколы, включая HTTP, Telnet, FTP, DNS и SSH).

На основании этих идентификаторов сетевой экран либо пропускает, либо отклоняет пакет передаваемых данных. Во втором случае отправитель может получить уведомление об ошибке.

\ManualSubsection{Типы межсетевых экранов}
\vspace*{-1.5em}
\ManualSubsubsection{Stateless FireWalls}

Работа файрвола основана на последовательном анализе сетевых пакетов и применении к ним набора правил. Исторически первым методом стала статическая пакетная фильтрация, которая анализирует только заголовки пакетов --- IP-адреса, порты и тип протокола. Такой подход прост, но не учитывает контекст соединения, что делает его уязвимым для спуфинга и сложных атак.

\ManualSubsubsection{Stateful FireWalls}

Более совершенным методом является фильтрация с учётом состояния. Они отличаются от статических фильтров способностью отслеживать и фильтровать трафик на основе данных о состоянии подключения, которые хранятся в таблице состояния. Чаще всего фильтрация осуществляется на основе правил, установленных администратором во время настройки компьютера и сетевого экрана. Однако таблица состояния позволяет файрволам с динамической фильтрацией самостоятельно принимать решения, используя данные о предыдущих подключениях для обучения. Например, типы трафика, которые провоцировали сбои в работе, будут отсеиваться.

\ManualSubsubsection{Proxy FireWalls}

Для защиты от угроз, скрытых в разрешённом трафике, применяется фильтрация на уровне приложений. Файрвол выступает в роли посредника (прокси), полностью «понимая» такие протоколы, как HTTP или FTP, и может блокировать отдельные вредоносные команды или контент внутри них.

\ManualSubsubsection{Next-Generation FireWalls}

Эволюция угроз привела к появлению файрволов нового поколения и глубокой инспекции пакетов (Deep Packet Inspection). Эти системы расширяют stateful-фильтрацию, анализируя само содержимое пакетов. Это позволяет реализовать тонкий контроль на уровне приложений, пользователей или категорий веб-ресурсов, а также выявлять угрозы, замаскированные под доверенный трафик.

\vspace{1.5em}

\noindent Таблица 5 – Сравнение типов межсетевых экранов

{\small
\noindent
\begin{tabular}{|p{0.20\textwidth}|p{0.38\textwidth}|p{0.36\textwidth}|}
\hline
\textbf{Тип} & \textbf{Характеристика} & \textbf{Применение} \\
\noalign{\hrule height 1.2pt}
\hline
Stateless & Анализ заголовков пакетов без контекста & Простые сети, базовая фильтрация \\
\hline
Stateful & Отслеживание состояния соединений & Корпоративные сети, средний уровень защиты \\
\hline
Proxy & Полное понимание протоколов приложений & Защита от угроз на уровне приложений \\
\hline
NGFW & Глубокая инспекция, контроль приложений & Современные корпоративные сети, высокий уровень защиты \\
\hline
\end{tabular}
}

\vspace{1em}

\ManualSubsection{Файрволы в современном мире}

Роль файрвола кардинально меняется в условиях современных трендов. Концепция Zero Trust («Никогда не доверяй, всегда проверяй») перемещает акцент с защиты только периметра на контроль доступа между любыми сегментами сети, где файрвол становится ключевым элементом микросетевой сегментации. \cite{kindervag2010zerotrust} описывает принципы Zero Trust как основу современной сетевой безопасности.

Распространение облачных технологий и удалённой работы породило модель FWaaS (Firewall-as-a-Service), которая позволяет применять единые политики безопасности ко всем пользователям и филиалам напрямую через облако, независимо от их физического местоположения. Современные NGFW всё чаще выступают как интегрированная платформа, автоматически получающая обновления об угрозах, передающая данные в системы SIEM и взаимодействующая с EDR для комплексного реагирования.

Современный файрвол эволюционировал из простого сетевого фильтра в интеллектуальную, интегрируемую платформу контроля доступа. Он остаётся краеугольным камнем сетевой безопасности, но его эффективность возможна только в рамках многоуровневой стратегии защиты (Defense-in-Depth). Надёжность обеспечивается сочетанием правильно настроенного файрвола с системами обнаружения вторжений, средствами защиты рабочих станций, регулярным обновлением ПО и шифрованием. Постоянный мониторинг и адаптация к меняющейся модели угроз превращают файрвол из статичного барьера в динамический инструмент управления сетевыми рисками.

\newpage
\refstepcounter{section}
\addcontentsline{toc}{section}{\protect\numberline{\thesection}Использование антивирусных средств и методов защиты от вредоносного ПО}
\begin{center}\textbf{\thesection\ ИСПОЛЬЗОВАНИЕ АНТИВИРУСНЫХ СРЕДСТВ И МЕТОДОВ ЗАЩИТЫ ОТ ВРЕДОНОСНОГО ПО}\end{center}
\vspace*{-1.5em}
\ManualSubsection{Введение в проблематику}

В современном цифровом мире проблема защиты информации вышла на первый план. Вредоносное программное обеспечение (или malware) стало основным инструментом киберпреступников, а ущерб от его деятельности исчисляется триллионами рублей ежегодно. От программ-вымогателей до шпионского ПО, крадущего персональные данные обычных пользователей, --- спектр угроз невероятно широк.

Использование антивирусных средств сегодня является не просто рекомендацией, а необходимостью. \cite{szor2005art} определяет антивирусную защиту как многоуровневый подход к обнаружению и нейтрализации вредоносного кода. Однако важно понимать, что современный антивирус --- это лишь один, хотя и важнейший, элемент в многоуровневой системе защиты.

\ManualSubsection{Классификация вредоносного ПО}

Вредоносное ПО представляет собой любую программу, созданную для нанесения ущерба компьютеру, сети или пользователю. Его классификация обширна, и каждая разновидность преследует свои цели.

Вирусы и черви являются классическими, самовоспроизводящимися угрозами. Вирус внедряется в легитимные файлы, а червь распространяется по сети самостоятельно, используя уязвимости. Троянские программы (трояны) маскируются под полезный софт, но выполняют скрытые вредоносные действия: крадут данные, открывают доступ к системе, загружают другие угрозы.

Программы-шпионы тайно собирают информацию: пароли, банковские реквизиты, историю браузера, переписку. Программы-вымогатели представляют собой одну из самых опасных современных угроз. Они шифруют файлы на диске и требуют выкуп за ключ расшифровки. Атакам подвергаются как частные лица, так и целые корпорации и государственные учреждения.

Криптоджекеры скрытно используют ресурсы компьютера (процессор, видеокарту) для добычи криптовалюты, что приводит к замедлению работы и износу оборудования. Рекламное ПО навязчиво показывает рекламу, часто перенаправляет поисковые запросы и может отслеживать поведение пользователя.

\vspace{1.5em}

\noindent Таблица 6 – Типы вредоносного программного обеспечения

{\small
\noindent
\begin{tabular}{|p{0.20\textwidth}|p{0.38\textwidth}|p{0.36\textwidth}|}
\hline
\textbf{Тип} & \textbf{Механизм действия} & \textbf{Последствия} \\
\noalign{\hrule height 1.2pt}
\hline
Вирусы & Внедрение в легитимные файлы & Повреждение файлов, распространение \\
\hline
Черви & Самостоятельное распространение по сети & Перегрузка сети, массовое заражение \\
\hline
Трояны & Маскировка под полезное ПО & Кража данных, удалённый доступ \\
\hline
Программы-вымогатели & Шифрование файлов & Потеря доступа к данным, требование выкупа \\
\hline
Шпионское ПО & Скрытый сбор информации & Утечка конфиденциальных данных \\
\hline
Криптоджекеры & Использование ресурсов для майнинга & Замедление работы, износ оборудования \\
\hline
\end{tabular}
}

\vspace{1em}

Цели атак разнообразны: от прямой финансовой выгоды (вымогательство, кража денег) до кибершпионажа, саботажа или создания сетей заражённых компьютеров (ботнетов) для дальнейших атак.

\ManualSubsection{Механизмы работы антивирусных средств}

Антивирусные движки используют многоуровневый подход, комбинируя несколько методов анализа.

Сигнатурный анализ представляет собой базовый метод. Каждый вредоносный файл имеет уникальный «цифровой отпечаток» (сигнатуру). Антивирус сверяет проверяемые файлы с базой известных сигнатур. Эффективен против массовых, известных угроз, но бесполезен против новых, неизвестных вирусов (так называемые угрозы «нулевого дня»).

Эвристический анализ является более продвинутым методом. Антивирус анализирует код программы на наличие подозрительных команд, структур и алгоритмов, характерных для вредоносного ПО. Это позволяет выявлять новые модификации известных вирусов или даже совершенно новые угрозы.

Поведенческий анализ представляет собой один из ключевых современных методов. Антивирус в реальном времени отслеживает поведение всех запущенных программ и процессов. Если приложение пытается выполнить подозрительные действия (например, массово зашифровать файлы, внести изменения в системный реестр или получить доступ к критическим областям памяти), оно будет заблокировано, даже если его сигнатура неизвестна.

Технология «песочницы» особенно важна для защиты от сложных угроз. Подозрительный файл запускается в полностью изолированной виртуальной среде («песочнице»), где анализируются все его действия. Если поведение признано вредоносным, файл блокируется, и реальная система остаётся нетронутой.

Облачный анализ и машинное обучение (ИИ) представляют собой будущее и настоящее антивирусных технологий. Подозрительные объекты могут мгновенно проверяться в облачных базах, где работают сложные алгоритмы искусственного интеллекта. Эти алгоритмы обучены на миллионах образцов и могут выявлять новые, ранее не встречавшиеся угрозы по косвенным признакам и аномальным шаблонам поведения.

\ManualSubsection{Комплексная стратегия защиты}

Установка даже самого совершенного антивируса не даёт 100\% гарантии. Безопасность обеспечивается только комплексным подходом, где технологические меры подкрепляются грамотными действиями пользователя.

Антивирус и файрвол работают совместно: антивирус защищает от внутренних угроз (файлов), а файрвол (брандмауэр) контролирует сетевой трафик, блокируя несанкционированные подключения извне и попытки передать данные из вашей системы.

Регулярное обновление критически важно для своевременной установки обновлений операционной системы, всех программ и самого антивируса. Большинство атак используют известные уязвимости, для которых уже выпущены «заплатки».

Осознанное поведение пользователя является самым важным слоем защиты. Не следует открывать вложения и переходить по ссылкам в подозрительных письмах, загружать программы только с официальных сайтов разработчиков, быть осторожным с флеш-накопителями от незнакомых источников и не использовать простые и повторяющиеся пароли.

Резервное копирование представляет собой главное оружие против программ-вымогателей. Нужно регулярно создавать резервные копии всех важных данных на внешний диск или в надёжное облачное хранилище. Правило «3-2-1»: три копии данных, на двух разных типах носителей, одна из которых хранится отдельно (оффлайн).

Использование сложных паролей и двухфакторной аутентификации (2FA) обеспечивает, что даже если пароль будет украден, доступ без второго фактора (код из SMS, приложения) будет невозможен. Работа под учётной записью с ограниченными правами не позволит многим видам вредоносного ПО внести критические изменения в систему.

\newpage
\refstepcounter{section}
\addcontentsline{toc}{section}{\protect\numberline{\thesection}Принципы резервного копирования и восстановления данных}
\begin{center}\textbf{\thesection\ ПРИНЦИПЫ РЕЗЕРВНОГО КОПИРОВАНИЯ И ВОССТАНОВЛЕНИЯ ДАННЫХ}\end{center}
\vspace*{-1.5em}
\ManualSubsection{Введение в резервное копирование}

Резервное копирование представляет собой процесс создания копий данных для защиты информации от потерь в случае повреждения, удаления или недоступности оригинальных данных. В современном цифровом мире, где данные являются одним из самых ценных активов организации, резервное копирование перестало быть опциональной мерой и стало критической необходимостью.

Потеря данных может произойти по множеству причин: от аппаратных сбоев жёстких дисков до кибератак с использованием вредоноса, от случайного удаления файлов до стихийных бедствий. \cite{preston2007backup} определяет резервное копирование как фундаментальный компонент стратегии восстановления после сбоев. Статистика показывает, что жёсткие диски имеют среднюю наработку на отказ (MTBF) в пределах 100--200 тысяч часов, что соответствует примерно 5--10 годам непрерывной работы. Однако отказ может произойти неожиданно в любой момент.

Восстановление данных с повреждённого оборудования является дорогостоящей и часто неполной процедурой, требующей помощи специалистов и специального оборудования. Наличие актуальной резервной копии позволяет восстановить работоспособность системы за часы, тогда как без бэкапа потеря может быть необратимой.

\ManualSubsection{Нормативные требования}

Необходимость проведения резервного копирования закреплена в международных стандартах и нормативных документах. Стандарт ISO/IEC 27001, являющийся мировым эталоном управления информационной безопасностью, содержит требование 8.13 о резервном копировании сведений, а также требование 8.14 об избыточности информации. \cite{iso27001} определяет требования к управлению информационной безопасностью. Эти требования обязывают организации разработать и внедрить политику резервного копирования, которая обеспечивает восстановление данных и систем в случае инцидентов.

В России требования к резервному копированию содержатся в Федеральном законе «О персональных данных», который предписывает организациям, обрабатывающим персональные данные, обеспечивать их сохранность и возможность восстановления. \cite{personaldata} устанавливает правовые основы защиты персональных данных. Организации, обрабатывающие персональные данные, должны регулярно проверять работоспособность резервных копий, иначе они рискуют получить штрафы от надзорных органов.

NIST (Национальный институт стандартов и технологий США) в своём документе NIST SP 800-184 «Guide for Cybersecurity Event Recovery» определяет резервное копирование как одну из ключевых составляющих плана восстановления после киберинцидентов.

\ManualSubsection{Правило 3-2-1}

Профессиональный стандарт в области резервного копирования сформулирован в виде правила 3-2-1, которое признано индустрией как оптимальная стратегия защиты данных. Это правило состоит из трёх компонентов.

Три копии данных означают, что у вас должно быть минимум три копии всех критичных данных --- один оригинал и две резервные копии. Наличие нескольких независимых копий обеспечивает множественность защиты: даже если одна или две копии будут повреждены или утеряны, у вас останется третья. Вероятность одновременного повреждения всех трёх копий катастрофически мала.

Два разных типа носителя означают, что эти копии должны храниться на двух разных типах носителей --- например, один на традиционном жёстком диске (HDD), другой на твёрдотельном накопителе (SSD). Разные типы носителей имеют разные механизмы отказа: HDD может заклинить из-за проблем с механикой, SSD может потерять электрический заряд. Использование различных технологий хранения делает менее вероятным, что все копии отказывают одновременно.

Одна копия в отдалённом месте означает, что как минимум одна копия должна храниться в физически отдалённом месте от основной инфраструктуры --- в другом здании, в другом городе или в облачном хранилище. Это защищает ваши данные от локальных катастроф, таких как пожары, наводнения, кражи или техногенные аварии, которые могут уничтожить локальную инфраструктуру.

Правило 3-2-1 создаёт многоуровневую защиту от различных типов инцидентов. Организация, следующая этому правилу, может практически гарантировать, что данные не будут полностью потеряны при любом разумном сценарии сбоя.

\ManualSubsection{Типы резервного копирования}

Существует несколько стратегий создания резервных копий, каждая с собственными характеристиками хранения, скорости создания и восстановления.

\ManualSubsubsection{Полное резервное копирование}

Полное резервное копирование представляет собой создание полной копии всех выбранных данных. При выполнении полного бэкапа копируется буквально каждый файл, каждая папка, все системные файлы, базы данных и конфигурации. Это самый простой и понятный тип резервного копирования, так как содержит всю необходимую информацию в одной копии.

Основное преимущество полного бэкапа состоит в том, что восстановление данных происходит максимально быстро и просто. Вам нужна только одна копия для полного восстановления. Кроме того, полные бэкапы не зависят от целостности предыдущих копий, что делает их надёжным фундаментом для системы резервного копирования.

Однако у полного бэкапа есть существенный недостаток: он занимает значительное место на хранилище и требует много времени для создания. Поэтому полные бэкапы обычно делаются нечасто --- например, раз в неделю или раз в месяц --- и служат основой для более быстрых дополнительных копий.

\ManualSubsubsection{Инкрементальное резервное копирование}

Инкрементальное резервное копирование использует совершенно другой подход к экономии времени и места. После создания одного полного бэкапа в каждый последующий момент копируются только те данные, которые изменились с момента последнего резервного копирования, независимо от того, был ли это полный бэкап или предыдущий инкрементальный.

Главное преимущество инкрементального резервного копирования --- это минимальное использование сетевой пропускной способности и пространства для хранения. Ежедневные инкрементальные копии могут быть на порядки меньше полного бэкапа.

Однако инкрементальное копирование имеет существенный недостаток при восстановлении: для полного восстановления данных вам нужны все копии по цепочке. Вам потребуется полный бэкап плюс все инкрементальные бэкапы в правильном порядке вплоть до нужного вам момента времени. Если в этой цепочке хотя бы одна копия повреждена или недоступна, восстановление полных данных становится невозможным.

\ManualSubsubsection{Дифференциальное резервное копирование}

Дифференциальное резервное копирование представляет собой компромисс между полным и инкрементальным методами. При дифференциальном копировании каждый раз копируются все данные, которые изменились с момента последнего полного бэкапа (не с момента последнего копирования вообще, а именно с момента последнего полного).

Получается, что дифференциальные копии растут со временем, так как накапливают всё больше изменений с момента последнего полного бэкапа. Однако преимущество дифференциального метода в том, что для восстановления нужны только два компонента: самый последний полный бэкап и последний дифференциальный бэкап.

Восстановление при дифференциальном методе происходит быстрее, чем при инкрементальном, потому что вам не нужно последовательно применять множество правок. Дифференциальное копирование обеспечивает хороший баланс между скоростью создания копий, объёмом хранилища и простотой восстановления.

\vspace{15em}

\noindent Таблица 7 – Сравнение типов резервного копирования

{\small
\noindent
\renewcommand{\arraystretch}{1.4}
\setlength{\tabcolsep}{3pt}
\begin{tabular}{|p{0.16\textwidth}|p{0.26\textwidth}|p{0.28\textwidth}|p{0.24\textwidth}|}
\hline
\textbf{Тип} & \textbf{Что копируется} & \textbf{Преимущества} & \textbf{Недостатки} \\
\noalign{\hrule height 1.2pt}
\hline
Полное & Все данные & Простое восстанов\-ление, независимость & Большой объём, долгое выполнение \\
\hline
Инкремент\-альное & Изменения с последнего копирования & Минимальный объём, быстрота & Сложное восстанов\-ление, зависимость от цепочки \\
\hline
Дифференц\-иальное & Изменения с последнего полного & Баланс объёма и скорости восстанов\-ления & Растущий объём копий \\
\hline
\end{tabular}
}

\vspace{1em}

\ManualSubsection{Метрики RTO и RPO}

При планировании системы резервного копирования критически важно понимать две ключевые метрики: Recovery Time Objective (RTO) и Recovery Point Objective (RPO).

RTO (Recovery Time Objective) представляет собой максимально допустимое время, в течение которого информационная система может быть недоступна. Например, если для веб-приложения электронной коммерции установлен RTO в 2 часа, это означает, что организация может позволить себе быть недоступной максимум 2 часа, иначе начнётся значительная потеря доходов. RTO определяет, насколько быстро должны происходить восстановление данных и запуск систем в рабочее состояние.

RPO (Recovery Point Objective) представляет собой максимально допустимое количество данных, которые организация готова потерять в результате инцидента. Обычно измеряется в единицах времени. Например, RPO в 1 час означает, что в худшем случае организация потеряет работу последнего часа. Если вы делаете резервные копии каждые 15 минут, ваше RPO составляет примерно 15 минут. RPO определяет, как часто вам необходимо выполнять резервное копирование.

Эти метрики напрямую влияют на выбор стратегии. Критичная база данных с RTO в 1 час и RPO в 15 минут требует частых бэкапов и быстрого восстановления. Архивные данные с RTO в 24 часа и RPO в 1 неделю требуют только еженедельных бэкапов.

\ManualSubsection{Варианты хранения резервных копий}

Локальное хранилище означает, что резервные копии хранятся на физических носителях, находящихся в пределах организации или на её близкой дистанции. Главное преимущество локального хранилища --- это скорость доступа. Восстановление данных с локального накопителя происходит с максимальной скоростью. Однако локальное хранилище подвержено локальным катастрофам.

Облачное хранилище использует удалённые серверы крупных поставщиков. Основное преимущество облака состоит в географической распределённости. Даже если в вашем городе произойдёт локальная катастрофа, ваши данные будут в безопасности на удалённых серверах. Недостатки облака включают текущие затраты на хранение и зависимость восстановления от скорости интернета.

Лучшая практика в индустрии --- использовать оба подхода одновременно, создавая гибридный подход: ежедневно создавать резервные копии на локальный NAS для быстрого восстановления, еженедельно загружать важные данные в облако для защиты от локальных катастроф, ежемесячно архивировать давние данные на ленту, которая хранится в отдельном здании.

\newpage
\section*{}
\vspace*{-2.5em}
\begin{center}\textbf{ЗАКЛЮЧЕНИЕ}\end{center}
\phantomsection
\addcontentsline{toc}{section}{Заключение}

Аутентификация и управление доступом в компьютерных сетях опираются на комплекс взаимосвязанных технологий и организационных мер. Корректная идентификация пользователей и устройств, проверка подлинности, назначение минимально необходимых прав и постоянный контроль их использования невозможны без понимания того, как устроена сама сетевая инфраструктура.

Базой для безопасного доступа служат принципы построения сетей и их архитектур: выбор топологий, знание уровневых моделей OSI и TCP/IP, понимание работы протоколов и сетевых сервисов. От того, насколько грамотно организованы локальные и глобальные сети, сегментация и маршрутизация трафика, зависят как эффективность работы приложений, так и возможности для разграничения доступа и обнаружения атак.

IP-адресация и маршрутизация являются техническим фундаментом для реализации политик доступа. Подсети, маски, маршруты и таблицы маршрутизации позволяют отделять одни домены доверия от других, реализовывать демилитаризованные зоны, граничные сегменты и изолированные контуры. Именно на этом уровне закладываются границы, внутри которых затем работают механизмы аутентификации и авторизации.

Межсетевые экраны и другие сетевые средства защиты выполняют роль фильтров, обеспечивающих прохождение только разрешённого трафика в соответствии с политиками безопасности. Они связывают уровни сети и прикладной логики, контролируя как адреса и порты, так и содержимое соединений. В связке с системами обнаружения и предотвращения вторжений это формирует внешний и внутренний периметр защиты.

Вредоносное программное обеспечение остаётся одним из основных инструментов компрометации учётных записей и обхода механизмов аутентификации. Поэтому антивирусная защита, системы контроля целостности, регулярное обновление программного обеспечения и обучение пользователей практике безопасной работы в сети являются неотъемлемой частью обеспечения доступа. Без защиты конечных точек даже самая совершенная сетевая архитектура остаётся уязвимой.

Резервное копирование и планирование восстановления завершают картину, обеспечивая аспект доступности в классической триаде конфиденциальность--целостность--доступность. Даже при успешной атаке или сбое инфраструктуры надлежащим образом организованные резервные копии и отработанные процедуры восстановления позволяют вернуть систему к рабочему состоянию и восстановить доступ к критичным данным в заданные сроки.

Таким образом, аутентификация и доступ не могут рассматриваться изолированно как набор логических проверок на уровне приложений. Это сквозная задача, затрагивающая архитектуру сетей, адресацию и маршрутизацию, межсетевые экраны, защиту конечных точек и механизмы резервного копирования. Только комплексный подход, сочетающий технические, организационные и процедурные меры, позволяет выстроить надёжную систему аутентификации и управления доступом в современных распределённых информационных системах.

\newpage
\begin{center}\textbf{СПИСОК СОКРАЩЕНИЙ}\end{center}
\phantomsection
\addcontentsline{toc}{section}{Список сокращений}

\vspace{1em}

{\small
\noindent
\begin{tabular}{|p{0.15\textwidth}|p{0.80\textwidth}|}
\hline
\textbf{Сокращение} & \multicolumn{1}{c|}{\textbf{Расшифровка}} \\
\noalign{\hrule height 1.2pt}
\hline
ACL & Access Control List (список контроля доступа) \\
\hline
AS & Autonomous System (автономная система) \\
\hline
BGP & Border Gateway Protocol (протокол граничного шлюза) \\
\hline
CIDR & Classless Inter-Domain Routing (бесклассовая междоменная маршрутизация) \\
\hline
DHCP & Dynamic Host Configuration Protocol (протокол динамической конфигурации хоста) \\
\hline
DNS & Domain Name System (система доменных имён) \\
\hline
EDR & Endpoint Detection and Response (обнаружение и реагирование на конечных точках) \\
\hline
HTTP & HyperText Transfer Protocol (протокол передачи гипертекста) \\
\hline
HTTPS & HyperText Transfer Protocol Secure (защищённый протокол передачи гипертекста) \\
\hline
IP & Internet Protocol (межсетевой протокол) \\
\hline
LAN & Local Area Network (локальная вычислительная сеть) \\
\hline
MAC & Media Access Control (управление доступом к среде) \\
\hline
NAT & Network Address Translation (преобразование сетевых адресов) \\
\hline
NGFW & Next-Generation Firewall (межсетевой экран нового поколения) \\
\hline
OSI & Open Systems Interconnection (взаимодействие открытых систем) \\
\hline
OSPF & Open Shortest Path First (сначала открыть кратчайший путь) \\
\hline
QoS & Quality of Service (качество обслуживания) \\
\hline
RPO & Recovery Point Objective (целевая точка восстановления) \\
\hline
RTO & Recovery Time Objective (целевое время восстановления) \\
\hline
SSH & Secure Shell (защищённая оболочка) \\
\hline
TCP & Transmission Control Protocol (протокол управления передачей) \\
\hline
TLS & Transport Layer Security (защита транспортного уровня) \\
\hline
UDP & User Datagram Protocol (протокол пользовательских дейтаграмм) \\
\hline
VLAN & Virtual Local Area Network (виртуальная локальная сеть) \\
\hline
VPN & Virtual Private Network (виртуальная частная сеть) \\
\hline
WAN & Wide Area Network (глобальная вычислительная сеть) \\
\hline
\end{tabular}
}

\clearpage
\nocite{*}
\printbibliography[heading=bibliography]

\end{document}
